% Guia de uso do template ppgcont.sty: mostra cada elemento formatado.
% Compilar: latexmk -pdf -outdir=build guia-do-template.tex
\documentclass[12pt,a4paper]{article}
\usepackage{ppgcont}
\addbibresource{guia.bib}

\titulo{Guia de uso do template do PPGCONT em \textit{LaTeX}}
\autor{Nome e Sobrenome}
\orientador{Orientador: Prof. Dr. Nome Sobrenome}
\ano{2027}

\begin{document}
\capa
\folhaderosto
\sumario

\section{Introdução}

O texto corrido sai em Times 12, justificado, com entrelinhas 1,5 e recuo de
1,25 cm na primeira linha. Termos estrangeiros vão em itálico, como
\textit{disclosure}. A citação entre parênteses usa \verb|\parencite|
\parencite{jensen1976,dickinson2011} e a citação no texto usa \verb|\textcite|:
\textcite{jensen1976} tratam dos custos de agência. Página específica:
\parencite[p.~310]{jensen1976}.

\subsection{Problema de pesquisa}

A seção secundária sai em maiúsculas, sem negrito. Basta usar o comando
\verb|\subsection| com o título em caixa normal; as maiúsculas são automáticas.

\subsection{Objetivos}

\subsubsection{Objetivo Geral}

A seção terciária sai em negrito, com a caixa que você digitar.

\subsubsection{Objetivos Específicos}

Os objetivos específicos entram em uma lista \verb|enumerate|:
\begin{enumerate}
  \item identificar as condições em que o primeiro item se aplica, com texto
        longo o bastante para mostrar o alinhamento da segunda linha;
  \item estimar o segundo item;
  \item comparar o terceiro item.
\end{enumerate}

\paragraph{Seção quaternária}

A seção quaternária (\verb|\paragraph|) sai sem destaque.

\section{Referencial teórico}

\subsection{Hipóteses de pesquisa}

A hipótese usa \verb|\hipotese{1}{texto}|:

\hipotese{1}{Há uma relação positiva entre a variável de interesse e a
variável dependente.}

\section{Procedimentos metodológicos}

A Tabela~\ref{tab:exemplo} mostra o padrão de tabela e a
Figura~\ref{fig:exemplo}, o de figura. A Equação~\ref{eq:ecl} é numerada à
direita.

\begin{table}[htbp]
  \caption{Variáveis dependentes}
  \label{tab:exemplo}
  \begin{tabularx}{\textwidth}{l l X l}
    \toprule
    \textbf{Variável} & \textbf{Sigla} & \textbf{Métrica} & \textbf{Suporte teórico} \\
    \midrule
    Erro de calibração & ErrCal & Diferença entre a PD estimada e a PD
      verdadeira do mecanismo gerador & \textcite{dickinson2011} \\
    Perda esperada & ECL & Produto de PD, LGD e EAD & \textcite{jensen1976} \\
    \bottomrule
  \end{tabularx}
  \nota{PD: probabilidade de inadimplência; LGD: perda dada a inadimplência;
    EAD: exposição no momento da inadimplência.}
  \fonte{Elaborada pelo autor (2027).}
\end{table}

\begin{equation}
  \mathrm{ECL}_i = \mathrm{PD}_i \times \mathrm{LGD}_i \times \mathrm{EAD}_i
  \label{eq:ecl}
\end{equation}

\begin{figure}[htbp]
  \centering
  \fbox{\parbox[c][3cm][c]{0.6\textwidth}{\centering (figura)}}
  \caption{Desenho de pesquisa}
  \label{fig:exemplo}
  \fonte{Elaborada pelo autor (2027).}
\end{figure}

\section{Cronograma}

\begin{table}[htbp]
  \caption{Cronograma}
  \label{tab:cronograma}
  \begin{tabularx}{\textwidth}{X *{6}{c}}
    \toprule
    \multirow{2}{*}{\textbf{Etapas}} & \multicolumn{6}{c}{\textbf{2027}} \\
    \cmidrule(l){2-7}
     & \textbf{Abr} & \textbf{Maio} & \textbf{Jun} & \textbf{Jul} & \textbf{Ago} & \textbf{Set} \\
    \midrule
    Ajustes no projeto        & X &   &   &   &   &   \\
    Execução dos experimentos &   & X & X & X &   &   \\
    Redação dos resultados    &   &   &   & X & X &   \\
    Defesa da dissertação     &   &   &   &   &   & X \\
    \bottomrule
  \end{tabularx}
  \fonte{Elaborada pelo autor (2027).}
\end{table}

\section{Resultados esperados}

As referências vêm logo depois desta seção, sem quebra de página, em APA,
alinhadas à esquerda e separadas por uma linha em branco.

\printbibliography
\end{document}
